\documentclass{article}
\usepackage{amsmath,amssymb,amsthm}
\usepackage{algorithm}
\usepackage[noend]{algpseudocode}
\usepackage{enumitem}
\usepackage{hyperref}
\usepackage{geometry}
\geometry{margin=1in}
\newtheorem{theorem}{Theorem}
\newtheorem{lemma}{Lemma}
\newtheorem{assumption}{Assumption}
\newcommand{\E}{\mathbb{E}}
\newcommand{\norm}[1]{\left\|#1\right\|}
\newcommand{\grad}[1]{\nabla #1}
\begin{document}

%%%%%%%%%%%%%%%%%%%%%%%%%%%%%%%%%%%%%%%%%%%%%%%%%%%%%%%%%%%%%%%%%%%%%%%%%%%%%
\section*{Algorithm Name}
\textbf{Randomized Block Coordinate Descent (RBCD)} for non‑convex smooth objectives.  
At each iteration a single block of coordinates is sampled at random and only
that block is updated by a (possibly constant) stepsize gradient step.

%%%%%%%%%%%%%%%%%%%%%%%%%%%%%%%%%%%%%%%%%%%%%%%%%%%%%%%%%%%%%%%%%%%%%%%%%%%%%
\section*{Mathematical Formulation}
Let $w\in\mathbb{R}^{d}$ be partitioned into $B$ disjoint blocks
\[
w = \bigl(w^{(1)},w^{(2)},\dots ,w^{(B)}\bigr),\qquad 
w^{(b)}\in\mathbb{R}^{d_{b}},\;\;\sum_{b=1}^{B}d_{b}=d .
\]

For a differentiable (possibly non‑convex) objective $f:\mathbb{R}^{d}\to\mathbb{R}$
define the block gradient
\[
\grad_{(b)} f(w) \;:=\; \frac{\partial f(w)}{\partial w^{(b)}}\in\mathbb{R}^{d_{b}} .
\]

At iteration $t$ a block index $i_{t}\in\{1,\dots ,B\}$ is drawn from a
distribution $p=(p_{1},\dots ,p_{B})$, $p_{b}>0$, $\sum_{b}p_{b}=1$.
With a stepsize $\eta_{t}>0$ the update reads
\[
\boxed{
\begin{aligned}
w_{t+1}^{(i_{t})} &= w_{t}^{(i_{t})} - \eta_{t}\,\grad_{(i_{t})} f(w_{t}),\\
w_{t+1}^{(b)}   &= w_{t}^{(b)}\qquad\forall\,b\neq i_{t}.
\end{aligned}}
\tag{RBCD}
\]

All other blocks remain unchanged.  The algorithm stops after $T$ iterations
or when a prescribed tolerance on the (expected) gradient norm is met.

%%%%%%%%%%%%%%%%%%%%%%%%%%%%%%%%%%%%%%%%%%%%%%%%%%%%%%%%%%%%%%%%%%%%%%%%%%%%%
\section*{Pseudocode}
\begin{algorithm}[H]
\caption{Randomized Block Coordinate Descent}
\begin{algorithmic}[1]
\State {\bf Input:} initial point $w_{0}$, stepsize schedule $\{\eta_{t}\}_{t\ge0}$,
probability vector $p$, maximum iterations $T$.
\For{$t=0$ \textbf{to} $T-1$}
    \State Sample block $i_{t}\sim p$.
    \State Compute block gradient $g_{t}= \grad_{(i_{t})}f(w_{t})$.
    \State $w_{t+1}^{(i_{t})}= w_{t}^{(i_{t})}-\eta_{t}\, g_{t}$.
    \State $w_{t+1}^{(b)} = w_{t}^{(b)}$ for all $b\neq i_{t}$.
\EndFor
\State {\bf Return} $w_{T}$.
\end{algorithmic}
\end{algorithm}

%%%%%%%%%%%%%%%%%%%%%%%%%%%%%%%%%%%%%%%%%%%%%%%%%%%%%%%%%%%%%%%%%%%%%%%%%%%%%
\section*{Dry Run Example}
Consider the one‑dimensional function $f(w)=(w-3)^{2}$, thus $B=1$,
$\grad f(w)=2(w-3)$.  Use $w_{0}=0$, constant stepsize $\eta=0.1$, and run three
iterations.

\begin{center}
\begin{tabular}{c|c|c|c}
$t$ & $w_{t}$ & $\grad f(w_{t})$ & $f(w_{t})$\\\hline
0 & $0$ & $2(0-3)=-6$ & $9$\\
1 & $0-0.1(-6)=0.6$ & $2(0.6-3)=-4.8$ & $(0.6-3)^{2}=5.76$\\
2 & $0.6-0.1(-4.8)=1.08$ & $2(1.08-3)=-3.84$ & $(1.08-3)^{2}=3.6864$\\
3 & $1.08-0.1(-3.84)=1.464$ & $2(1.464-3)=-3.072$ & $(1.464-3)^{2}=2.3660$
\end{tabular}
\end{center}
The objective value decreases monotonically, illustrating the mechanics of
(RBCD) in the simplest setting.

%%%%%%%%%%%%%%%%%%%%%%%%%%%%%%%%%%%%%%%%%%%%%%%%%%%%%%%%%%%%%%%%%%%%%%%%%%%%%
\section*{Assumptions}
\begin{assumption}[Blockwise Smoothness]\label{ass:smooth}
For each block $b\in\{1,\dots ,B\}$ there exists $L_{b}>0$ such that for all
$w,\tilde w$ differing only in block $b$,
\[
f(\tilde w)\le f(w)+\langle \grad_{(b)}f(w),\tilde w^{(b)}-w^{(b)}\rangle
+\frac{L_{b}}{2}\norm{\tilde w^{(b)}-w^{(b)}}^{2}.
\tag{A1}
\]
Define $L_{\max}:=\max_{b}L_{b}$ and $L_{\mathrm{avg}}:=\sum_{b}p_{b}L_{b}$.
\end{assumption}

\begin{assumption}[Bounded Gradient]\label{ass:bound}
There exists $G>0$ such that for all $w$ and all blocks $b$,
\[
\norm{\grad_{(b)}f(w)}\le G . \tag{A2}
\]
\end{assumption}

\begin{assumption}[Stepsize]\label{ass:stepsize}
The stepsize sequence satisfies $0<\eta_{t}\le \frac{1}{L_{\max}}$
for all $t$.  In the constant‑stepsize case we write $\eta_{t}\equiv\eta$.
\end{assumption}

No convexity is required; $f$ may be non‑convex.

%%%%%%%%%%%%%%%%%%%%%%%%%%%%%%%%%%%%%%%%%%%%%%%%%%%%%%%%%%%%%%%%%%%%%%%%%%%%%
\section*{Main Convergence Theorem}
\begin{theorem}[Expected Gradient Norm Decay]\label{thm:main}
Let $\{w_{t}\}_{t\ge0}$ be generated by (RBCD) under Assumptions
\ref{ass:smooth}–\ref{ass:stepsize}.
Define the average iterate
\[
\bar w_{T}:=\frac{1}{T}\sum_{t=0}^{T-1} w_{t}.
\]
Then for any $T\ge1$,
\[
\frac{1}{T}\sum_{t=0}^{T-1}\E\!\left[
\sum_{b=1}^{B}p_{b}\,\norm{\grad_{(b)}f(w_{t})}^{2}\right]
\;\le\;
\frac{2\bigl(f(w_{0})-f^{\inf}\bigr)}{\eta\,T}
\;+\;
\eta\,L_{\mathrm{avg}}\,G^{2},
\tag{T1}
\]
where $f^{\inf}:=\inf_{w}f(w)$ and $\eta$ is any constant stepsize satisfying
$0<\eta\le 1/L_{\max}$.
Consequently, choosing $\eta = \Theta\!\bigl(1/\sqrt{T}\bigr)$ yields
\[
\min_{0\le t<T}\E\!\left[
\sum_{b}p_{b}\,\norm{\grad_{(b)}f(w_{t})}^{2}\right]
= O\!\left(\frac{1}{\sqrt{T}}\right).
\]
\end{theorem}

%%%%%%%%%%%%%%%%%%%%%%%%%%%%%%%%%%%%%%%%%%%%%%%%%%%%%%%%%%%%%%%%%%%%%%%%%%%%%
\section*{Theoretical Properties}
\begin{enumerate}[label={\arabic*.}]
\item \textbf{Stability.}
The bound (T1) guarantees that the expected squared block‑gradient is
summable; thus the iterates cannot diverge as long as the stepsize respects
Assumption~\ref{ass:stepsize}.

\item \textbf{Hyperparameter Sensitivity.}
The constant $L_{\mathrm{avg}}$ appears only linearly in the second term of
(T1).  Over‑estimating $L_{\max}$ merely forces a smaller admissible stepsize,
which slows convergence but does not break the guarantee.

\item \textbf{Comparison to Full‑Gradient Descent.}
If $B=1$ (single block) the bound reduces to the classical non‑convex SGD
rate $O(1/\sqrt{T})$ with variance term replaced by $L_{\max}G^{2}$.  For
$B>1$, the per‑iteration cost is reduced proportionally to the average block
size while the convergence rate in terms of \emph{iterations} remains the same.

\item \textbf{Special Cases.}
\begin{itemize}
\item \emph{Uniform sampling} $p_{b}=1/B$ yields $L_{\mathrm{avg}}=\frac{1}{B}\sum_{b}L_{b}$.
\item \emph{Importance sampling} $p_{b}\propto L_{b}$ minimizes the term
$\eta L_{\mathrm{avg}}$ and hence improves the constant in (T1).
\end{itemize}
\end{enumerate}

\section*{Discussion}
The theorem shows that even without convexity, a simple random block
selection combined with a fixed stepsize delivers a sublinear decay of the
expected squared block‑gradient.  The result matches the best known rates for
non‑convex stochastic gradient methods while exploiting the cheaper
per‑iteration cost of coordinate updates.  Extensions to mini‑batch blocks,
accelerated stepsizes, or variance‑reduced gradients follow the same
proof skeleton.

%%%%%%%%%%%%%%%%%%%%%%%%%%%%%%%%%%%%%%%%%%%%%%%%%%%%%%%%%%%%%%%%%%%%%%%%%%%%%
\section*{Preliminaries (Definitions and Lemmas)}
\begin{lemma}[One‑step Expected Descent]\label{lem:descent}
Under Assumptions~\ref{ass:smooth}–\ref{ass:stepsize},
for any iteration $t$,
\[
\E_{i_{t}}\!\bigl[f(w_{t+1})\mid w_{t}\bigr]
\le f(w_{t})-\frac{\eta}{2}\sum_{b=1}^{B}p_{b}\,
\norm{\grad_{(b)}f(w_{t})}^{2}
+\frac{\eta^{2}L_{\mathrm{avg}}}{2}\,G^{2}.
\tag{L1}
\]
\end{lemma}
\begin{proof}
Fix $w_{t}$ and denote $i=i_{t}$.  By blockwise smoothness (A1) with
$\tilde w = w_{t} -\eta\,\grad_{(i)}f(w_{t})$ we have
\[
f(w_{t+1})\le
f(w_{t})-\eta\!\langle \grad_{(i)}f(w_{t}),\grad_{(i)}f(w_{t})\rangle
+\frac{L_{i}\eta^{2}}{2}\norm{\grad_{(i)}f(w_{t})}^{2}.
\tag{S1}
\]
Since $\langle a,a\rangle =\norm{a}^{2}$, (S1) becomes
\[
f(w_{t+1})\le
f(w_{t})-\eta\bigl(1-\tfrac{L_{i}\eta}{2}\bigr)
\norm{\grad_{(i)}f(w_{t})}^{2}.
\tag{S2}
\]
Because $\eta\le 1/L_{\max}\le 2/L_{i}$, the factor in parentheses satisfies
$1-\tfrac{L_{i}\eta}{2}\ge \tfrac12$.  Hence
\[
f(w_{t+1})\le
f(w_{t})-\frac{\eta}{2}\norm{\grad_{(i)}f(w_{t})}^{2}.
\tag{S3}
\]
Taking expectation over $i$ conditioned on $w_{t}$,
\[
\E_{i}\!\bigl[f(w_{t+1})\mid w_{t}\bigr]
\le f(w_{t})-\frac{\eta}{2}\sum_{b=1}^{B}p_{b}\,
\norm{\grad_{(b)}f(w_{t})}^{2}.
\tag{S4}
\]
To incorporate the bounded‑gradient term we add and subtract
$\frac{\eta^{2}L_{i}}{2}\norm{\grad_{(i)}f(w_{t})}^{2}$ inside the expectation:
\[
\E_{i}\!\bigl[f(w_{t+1})\mid w_{t}\bigr]
\le f(w_{t})-\frac{\eta}{2}\sum_{b}p_{b}\norm{\grad_{(b)}f(w_{t})}^{2}
+\frac{\eta^{2}}{2}\sum_{b}p_{b}L_{b}\norm{\grad_{(b)}f(w_{t})}^{2}.
\tag{S5}
\]
Using $\norm{\grad_{(b)}f(w_{t})}\le G$ (Assumption~\ref{ass:bound}) and the
definition of $L_{\mathrm{avg}}$, the last term is bounded by
$\frac{\eta^{2}L_{\mathrm{avg}}}{2}G^{2}$, which yields (L1).  ∎
\end{proof}

%%%%%%%%%%%%%%%%%%%%%%%%%%%%%%%%%%%%%%%%%%%%%%%%%%%%%%%%%%%%%%%%%%%%%%%%%%%%%
\section*{Main Proof (Step‑by‑Step Derivation)}
We prove Theorem~\ref{thm:main} by telescoping the descent inequality of
Lemma~\ref{lem:descent}.

\begin{enumerate}
\item[Step 1.]  Apply Lemma~\ref{lem:descent} and take total expectation:
\[
\E\!\bigl[f(w_{t+1})\bigr]
\le \E\!\bigl[f(w_{t})\bigr]
-\frac{\eta}{2}\E\!\Bigl[\sum_{b}p_{b}\norm{\grad_{(b)}f(w_{t})}^{2}\Bigr]
+\frac{\eta^{2}L_{\mathrm{avg}}}{2}G^{2}.
\tag{P1}
\]

\item[Step 2.]  Rearrange (P1) to isolate the gradient term:
\[
\frac{\eta}{2}\E\!\Bigl[\sum_{b}p_{b}\norm{\grad_{(b)}f(w_{t})}^{2}\Bigr]
\le \E\!\bigl[f(w_{t})\bigr]-\E\!\bigl[f(w_{t+1})\bigr]
+\frac{\eta^{2}L_{\mathrm{avg}}}{2}G^{2}.
\tag{P2}
\]

\item[Step 3.]  Sum (P2) over $t=0,\dots ,T-1$:
\[
\frac{\eta}{2}\sum_{t=0}^{T-1}
\E\!\Bigl[\sum_{b}p_{b}\norm{\grad_{(b)}f(w_{t})}^{2}\Bigr]
\le \E\!\bigl[f(w_{0})\bigr]-\E\!\bigl[f(w_{T})\bigr]
+ \frac{\eta^{2}L_{\mathrm{avg}}}{2}G^{2}T.
\tag{P3}
\]

\item[Step 4.]  Since $f(w_{T})\ge f^{\inf}$, replace $-\E[f(w_{T})]$ by
$-f^{\inf}$ to obtain an upper bound:
\[
\frac{\eta}{2}\sum_{t=0}^{T-1}
\E\!\Bigl[\sum_{b}p_{b}\norm{\grad_{(b)}f(w_{t})}^{2}\Bigr]
\le f(w_{0})-f^{\inf}
+ \frac{\eta^{2}L_{\mathrm{avg}}}{2}G^{2}T.
\tag{P4}
\]

\item[Step 5.]  Divide both sides of (P4) by $\eta T$:
\[
\frac{1}{T}\sum_{t=0}^{T-1}
\E\!\Bigl[\sum_{b}p_{b}\norm{\grad_{(b)}f(w_{t})}^{2}\Bigr]
\le
\frac{2\bigl(f(w_{0})-f^{\inf}\bigr)}{\eta T}
+\eta L_{\mathrm{avg}} G^{2}.
\tag{P5}
\]

\item[Step 6.]  Inequality (P5) is exactly the bound (T1) claimed in
Theorem~\ref{thm:main}.  ∎
\end{enumerate}

%%%%%%%%%%%%%%%%%%%%%%%%%%%%%%%%%%%%%%%%%%%%%%%%%%%%%%%%%%%%%%%%%%%%%%%%%%%%%
\section*{Rate Derivation (With Constants)}
From (T1) we have for a constant stepsize $\eta$,
\[
\min_{0\le t<T}
\E\!\Bigl[\sum_{b}p_{b}\norm{\grad_{(b)}f(w_{t})}^{2}\Bigr]
\le
\frac{2\bigl(f(w_{0})-f^{\inf}\bigr)}{\eta T}
+\eta L_{\mathrm{avg}} G^{2}.
\]
Choosing $\eta = \sqrt{\dfrac{2\bigl(f(w_{0})-f^{\inf}\bigr)}{L_{\mathrm{avg}}G^{2}}\,\dfrac{1}{\sqrt{T}}$
balances the two terms and yields
\[
\min_{t<T}\E\!\Bigl[\sum_{b}p_{b}\norm{\grad_{(b)}f(w_{t})}^{2}\Bigr]
\le
2\sqrt{2\,L_{\mathrm{avg}}\,G^{2}\,\bigl(f(w_{0})-f^{\inf}\bigr)}\;\frac{1}{\sqrt{T}}
\;=\;O\!\left(\frac{1}{\sqrt{T}}\right).
\]

%%%%%%%%%%%%%%%%%%%%%%%%%%%%%%%%%%%%%%%%%%%%%%%%%%%%%%%%%%%%%%%%%%%%%%%%%%%%%
\section*{Special Cases}
\begin{itemize}
\item \textbf{Uniform sampling.} $p_{b}=1/B$ gives $L_{\mathrm{avg}}=\frac{1}{B}\sum_{b}L_{b}$.
\item \textbf{Deterministic cyclic block selection.} The same bound holds
with $p_{b}=1/B$ after a full cycle, because each block is visited exactly
once per $B$ iterations.
\item \textbf{Mini‑batch blocks.} If a set $\mathcal{S}_{t}$ of $m$ blocks is
sampled uniformly, the analysis proceeds with $\eta\le 1/(\max_{b\in\mathcal{S}_{t}}L_{b})$
and $p_{b}=m/B$, leading to a factor $m$ improvement in the descent term.
\end{itemize}

%%%%%%%%%%%%%%%%%%%%%%%%%%%%%%%%%%%%%%%%%%%%%%%%%%%%%%%%%%%%%%%%%%%%%%%%%%%%%
\end{document}